\paragraph*{Three models.} Three models score each scan for progression: an imaging model, a clinical-history model
and a combined model. The imaging model uses the 27 measurements and is given no dates. Its score therefore depends on
the tumour alone and can be set against the clinical-history score. The clinical-history model uses the treatment
course and the patient's characteristics and sees no images. The combined model uses both sets of inputs. PRISM
reports all three scores, so a reader can see whether a scan's combined score follows the images or the clinical
history.

Each model is a logistic regression with a ridge penalty, fitted with scikit-learn \citep{pedregosa2011}. Logistic
regression gives each input a weight and turns the weighted sum into a probability of progression. The ridge (L2)
penalty pulls the weights towards zero, which keeps them stable when there are many correlated inputs and few
patients. Deep networks fine-tuned on MU did not add to a ridge model of the clinical history and image measurements
(Table~\ref{tab:2}). Before fitting, each missing input is set to its median in the training data, an extra input
marks the scans where it was missing, and every input is standardised to mean zero and unit variance. The penalty
strength is set by its inverse, $C$, where smaller values pull harder. $C$ is chosen from six values (0.01, 0.03, 0.1,
0.3, 1 and 3) by threefold cross validation within the training data, with each patient's scans kept in one fold.
The chosen value is the one with the lowest mean log loss, which measures how much probability the model gives to the
true outcome. The model is then refitted with that value on all training scans.

The full clinical-history model cannot be computed at the two external centres, because neither records the full
treatment course. LUMIERE dates each scan only by the weeks since surgery. UCSF-ALPTDG records the days since
diagnosis, which it defines as the first surgery, and the start of radiotherapy, but not its end. The models are
applied there without retraining, so PRISM uses reduced models: they are fitted on MU, but only with the inputs that
all three cohorts record. The reduced clinical-history model takes age, sex, IDH and MGMT status, the days since
surgery, and a flag for scans in the first 12 weeks after radiotherapy. On MU the flag is computed from the end of
radiotherapy. At LUMIERE it is the rater's mark for scans taken less than three months after radiotherapy. At
UCSF-ALPTDG it marks scans 42 to 125 days after the start of radiotherapy, which assumes a six-week course. The
reduced combined model adds the log volumes of the enhancing tumour, of necrosis and cavity together, and of the FLAIR
hyperintensity. It also adds the change in each volume since the previous scan, whether there was a previous scan, and
the interval since it.

At UCSF-ALPTDG the supplied masks also mark the resection cavity, so PRISM applies the full imaging model. The one
earlier scan in the dataset serves as the patient's history. LUMIERE's supplied masks have three compartments, with
necrosis and cavity combined. Its imaging model is therefore fitted on MU with those three volumes and their change
since the previous scan. At the 10 LUMIERE scans with no previous scan, a model of the three current volumes is used
instead. The imaging model needs only the scans and their masks, so PRISM gives an imaging score at every centre,
including one that records no treatment dates.

% Sources: prism/audit/conformal.py; METHODS_PLAN.md M2; RELIABILITY_CARD.md; REVISION_DATA.md R2; refs
% angelopoulos2023, vovk2022book, romano2020, barber2021limits, johansson2015, ding2023, mehrtens2023, mitchell2019.
